% TT75040C
% Stromunfälle
% 14.0
% 2013-11-30
{\label{m:stromunfall}}
{
\begin{itemize}
  \item \EE{Eigenschutz!} Bevor der Patienten berührt wird,
  muss man sich vergewissern, dass der Patient
  keinen Kontakt mehr mit dem Stromkreis hat! \EE{Strom
  abschalten (lassen)!}
  \begin{itemize}
    \item Gefahrenbreich beachten:
    \prettyref{AASS-topic:gefahrenbereich-strom}
    \item Sicherheitsabstand
    (\prettyref{AASS-topic:gefahrenbereich-strom})
    \item Schrittspannung und Spannungskegel
    beachten (\prettyref{AASS-topic:gefahrenbereich-strom})
    \item Bahnanlagen sperren/sichern lassen
    (\prettyref{AASS-topic:gefahrenbereich-strom})
    \item Schalter ausschalten, (Haupt-)Sicherung
    auslösen, Starkstrom/Hochspannung durch Fachmann  abschalten bzw.
    erden lassen
    \item Strom gegen Wiedereinschalten sichern!
    \item Patient mit isolierendem Material aus dem Stromkreis befreien

  \end{itemize}
  \item Beurteilung der vitalen Bedrohung.
  \par \TxBeiVit
  \TxMassVitMK
  \item \EE{Monitoring und
  Reanimationsbereitschaft}: auf evtl.
  Herzrhythmusstörungen vorbereitet sein
  (bis zu 24\,h Latenz)
  \item Genaue Patientenuntersuchung ${\rightarrow}$
  Eintritts- und Austrittsmarke
  \item Hospitalisierung auch bei symptomlosen
  Patienten zur Überwachung
  \item Stromverbrennungen wie Verbrennungen versorgen
\end{itemize}
}
